\chapter{Breast Biopsy}

\section*{Overview}
Breast biopsy is the removal of a small sample of breast tissue for microscopic examination to determine whether a lump or suspicious imaging finding is benign or malignant. It is the definitive diagnostic test for breast lesions.
\section*{Alternative Names}
Core needle biopsy; vacuum-assisted biopsy; stereotactic breast biopsy; ultrasound-guided breast biopsy; excisional breast biopsy
\section*{Principle}
A needle or surgical instrument obtains breast tissue from a suspicious area, which is then examined by a pathologist. Sampling may be guided by ultrasound, mammography (stereotaxis), or MRI to target the abnormality precisely, and the presence of cancer cells in the sample establishes the diagnosis.
\section*{History}
Breast biopsy developed with the rise of surgical pathology in the late 19th and early 20th centuries, and image-guided core needle biopsy became the standard approach in the 1990s, replacing most open surgical biopsies.
\section*{Why It Is Done}
Ordered for a palpable breast mass, suspicious calcifications or masses on mammography, abnormal ultrasound or MRI findings, nipple discharge, or skin changes that raise concern for malignancy.
\section*{Is This a Primary or Secondary Test or Procedure}
Primary diagnostic procedure for breast lesions and essential to confirming or excluding breast cancer.
\section*{How It Is Performed}
The lesion is localized by palpation or imaging. After local anesthesia, a core needle or vacuum-assisted device is inserted through a small skin nick and multiple samples are taken from the lesion, often under ultrasound or stereotactic guidance. A small metal clip may be placed to mark the biopsy site.
\section*{Preparation}
Informed consent, review of anticoagulant use, and avoidance of aspirin or blood thinners before the procedure. No fasting is required for core needle biopsy.
\section*{Contraindications and Precautions}
Precautions include bleeding disorders or anticoagulation, and the procedure is deferred in patients with active breast infection. The patient should be able to remain still during image-guided sampling.
\section*{Risks}
Pain, bruising and hematoma, bleeding, infection, and rarely pneumothorax when sampling near the chest wall. Core needle biopsy has a small rate of false-negative results.
\section*{Aftercare and Recovery}
The site is bandaged and the patient applies ice and avoids heavy lifting for a day or two. The specimens are sent for histopathologic analysis, with results typically available within several days.
\section*{Outcome and Follow-up}
The pathology report classifies the lesion as benign, atypical, or malignant, and reports tumor type, grade, and hormone receptor status when cancer is found. A benign result that does not match the imaging appearance may require repeat biopsy or surgical excision.
\section*{Expected Results}
Benign findings such as fibroadenoma, fibrocystic change, or benign ductal tissue with no malignant cells.
\section*{Follow-up}
Patients with benign results may require imaging follow-up, while those with cancer proceed to staging and treatment planning, including surgery, radiation, and systemic therapy.
\section*{References}
\begin{enumerate}
  \item MedlinePlus. Breast biopsy. U.S. National Library of Medicine; 2025.
  \item Current Medical Diagnosis and Treatment. McGraw-Hill; 2025.
\end{enumerate}

\clearpage
